\documentclass[11pt]{article}

\usepackage[margin=1in]{geometry}
\usepackage[T1]{fontenc}
\usepackage[utf8]{inputenc}
\usepackage{lmodern}
\usepackage{microtype}
\usepackage{setspace}
\usepackage{amsmath, amssymb, amsfonts}
\usepackage{siunitx}
\usepackage{booktabs}
\usepackage{array}
\usepackage{longtable}
\usepackage{graphicx}
\usepackage{xcolor}
\usepackage{float}
\usepackage{hyperref}
\usepackage{enumitem}
\usepackage{listings}

\hypersetup{
    colorlinks=true,
    linkcolor=blue!50!black,
    urlcolor=blue!50!black,
    citecolor=blue!50!black
}

\onehalfspacing
\setlength{\parskip}{0.5em}
\setlength{\parindent}{0pt}

\lstdefinestyle{projectcode}{
    basicstyle=\ttfamily\small,
    keywordstyle=\color{blue!60!black}\bfseries,
    commentstyle=\color{green!35!black},
    stringstyle=\color{red!50!black},
    numbers=left,
    numberstyle=\tiny\color{gray},
    stepnumber=1,
    numbersep=8pt,
    showstringspaces=false,
    breaklines=true,
    frame=single,
    rulecolor=\color{black!20},
    tabsize=2
}

\newcommand{\placeholderfigure}[4]{
\begin{figure}[H]
\centering
\fbox{
    \begin{minipage}[c][#2][c]{0.9\linewidth}
    \centering
    \textbf{#1}\\[0.75em]
    \small #3
    \end{minipage}
}
\caption{#4}
\end{figure}
}

\title{\textbf{Technical Report:}\\
Hydrogen Fuel-Cell Burn-and-Coast Race Strategy Modeling,\\
Track-Aware Simulation, and Design Sensitivity Analysis}
\author{Race Strategy Capstone Project}
\date{\today}

\begin{document}
\maketitle

\begin{abstract}
This report documents a track-aware race strategy framework for a hydrogen fuel-cell competition vehicle operating under endurance-style constraints such as those encountered in Shell Eco-marathon. The project combines physics-based longitudinal modeling, curvature- and grade-aware track processing, a rule-based burn-and-coast controller, an offline Python simulation engine, and an interactive browser dashboard for rapid engineering studies. The resulting tool supports two closely related decisions: first, race-day strategy synthesis by identifying where the vehicle should accelerate, coast, and preserve time margin; and second, vehicle-design evaluation by exposing the sensitivity of energy consumption and event completion to mass, aerodynamic drag area, rolling resistance, traction limits, and cornering capability. The report presents the mathematical formulation, software architecture, competition relevance, example results, and implementation details, and includes appendices with representative code excerpts from the working system.
\end{abstract}

\section{Introduction and Competition Context}

Hydrogen fuel-cell prototype vehicles compete under a distinct engineering objective: travel a prescribed course while satisfying event timing rules and minimizing energy consumption. In a competition such as Shell Eco-marathon, a vehicle cannot simply move as slowly as possible to save energy; it must still finish within a timing window or maintain a required average pace. As a result, performance is governed by the interaction between vehicle design, track geometry, and strategy.

For the present case study, the relevant event requirement is four laps in \SI{35}{min}. Using the currently processed \texttt{SEM 2023 US} loop length of approximately \SI{3850}{m}, this corresponds to an event distance of \SI{15400}{m} and a required average speed of approximately \SI{26.4}{km/h}. That timing constraint is important because it prevents trivial low-speed energy minimization and makes strategy synthesis a true constrained design problem.

This interaction is especially important for vehicles with limited continuous propulsion power, low mass, and strong sensitivity to aerodynamic and rolling losses. A strategy that is acceptable on a flat oval may be suboptimal or even infeasible on a circuit with tighter curvature, elevation variation, or repeated acceleration zones. Conversely, a vehicle-design change that appears small in isolation, such as a modest reduction in drag area or rolling resistance, can materially alter the optimal burn-and-coast schedule across an entire event.

The present project addresses this problem by building a unified workflow that:

\begin{itemize}[leftmargin=1.5em]
\item converts measured GPS track data into simulation-ready geometry,
\item computes grade and curvature along the track,
\item simulates a hydrogen fuel-cell vehicle with propulsion and cornering constraints,
\item treats the horn, live telemetry system, and driver strategy dashboard as accessories-battery loads rather than fuel-cell loads,
\item applies a rule-based burn-and-coast controller with lookahead, and
\item exposes the results through an interactive engineering dashboard.
\end{itemize}

The value of this approach is practical as well as academic. It provides a race strategy tool that can inform driver instructions and system operating targets, while also acting as a design-analysis platform that shows how parameter changes propagate into competition outcomes.

\placeholderfigure
{Figure Placeholder 1: System-Level Workflow}
{2.4in}
{Recommended content: a block diagram showing raw CSV track data feeding the track exporter, then the simulation engine, then outputs and the web dashboard. Include vehicle parameters, strategy parameters, the traction electrical path, and a separate accessories battery branch for the horn, live telemetry, and driver dashboard.}
{System architecture of the race strategy framework. The figure should illustrate how raw measured track data, vehicle properties, and controller settings are transformed into simulation outputs such as speed traces, fuel-cell power demand, cumulative traction energy, and track maps colored by speed. It should also show that the horn, live telemetry system, and driver dashboard are powered by a separate accessories battery rather than the fuel-cell traction path.}

\section{Problem Definition and Objectives}

The project solves a constrained minimum-energy traversal problem. Given a track of known distance, elevation, and curvature, the objective is to complete the event within the available time while minimizing electrical energy consumption and, by extension, hydrogen use. The problem is subject to:

\begin{itemize}[leftmargin=1.5em]
\item aerodynamic, rolling, and grade-induced resistance,
\item wheel-force and motor-power limits,
\item fuel-cell usable power limits,
\item curvature-based cornering speed limits,
\item a maximum allowable speed cap, and
\item strategy logic that determines when to apply power and when to coast.
\end{itemize}

From an engineering standpoint, the project pursues four goals:

\begin{enumerate}[leftmargin=1.5em]
\item develop a physically interpretable simulation model suitable for a capstone-level technical study;
\item support real measured tracks rather than synthetic geometry alone;
\item provide a user-facing interface that enables rapid scenario exploration without editing source code; and
\item expose parameter dependencies so that the same framework can support both strategy optimization and vehicle-design decision-making.
\end{enumerate}

\section{Vehicle and Track Modeling}

\subsection{Track Representation}

The track is represented in the space domain using cumulative distance \(s\), with accompanying arrays for elevation, grade, curvature, latitude, and longitude. A raw track CSV contains latitude, longitude, and optionally elevation. The exporter resamples the raw path onto a uniform spatial grid, enabling stable numerical calculations and direct comparison of track segments.

For the current measured example, the processed \texttt{SEM 2023 US} track has a length of approximately \SI{3850}{m} and \(1926\) resampled points, corresponding to a grid spacing of roughly \SI{2}{m}. Uniform spatial spacing is useful because it aligns track features, controller lookahead, and engineering plots in terms of distance rather than arbitrary sample index.

\subsection{Vehicle Parameters}

The current baseline vehicle model tracks the quantities listed in Table~\ref{tab:parameters}. These parameters are intentionally exposed because they are precisely the quantities that designers and race teams adjust during iterative development.

\begin{table}[H]
\centering
\caption{Primary design and strategy parameters tracked by the framework.}
\label{tab:parameters}
\begin{tabular}{p{0.34\linewidth} p{0.2\linewidth} p{0.36\linewidth}}
\toprule
\textbf{Parameter} & \textbf{Representative Value} & \textbf{Engineering Significance} \\
\midrule
Vehicle mass \(m\) & \SI{185}{kg} & Governs inertial demand, grade losses, and rolling losses \\
Drag area \(C_dA\) & \SI{0.19}{m^2} & Dominant source of power demand at higher speed \\
Rolling resistance \(C_{rr}\) & 0.0045 & Strongly influences low-speed endurance performance \\
Drive efficiency \(\eta_{\mathrm{drive}}\) & 0.90 & Converts mechanical wheel demand into traction electrical demand \\
Motor power limit & \SI{1000}{W} & Caps propulsion power at moderate to high speed \\
Motor force limit & \SI{190}{N} & Caps low-speed tractive force and therefore launch/exit behavior \\
Lateral acceleration limit & \SI{1.40}{m/s^2} & Converts track curvature into allowable cornering speed \\
Fuel-cell peak usable power & \SI{850}{W} & Limits sustainable energy delivery to the traction system \\
Burn/coast thresholds & \SI{7.2}{m/s} / \SI{9.2}{m/s} & Control how aggressively the strategy preserves time margin \\
\bottomrule
\end{tabular}
\end{table}

\section{Mathematical Formulation}

\subsection{Longitudinal Force Balance}

The model is most naturally expressed in the space domain, because race strategy is evaluated by track position rather than by a preselected time horizon. The vehicle experiences tractive effort opposed by aerodynamic drag, rolling resistance, and grade force. The resistive forces are

\begin{align}
F_{\mathrm{aero}} &= \frac{1}{2}\rho C_dA v^2 \\
F_{\mathrm{roll}} &= mgC_{rr}\cos\theta \\
F_{\mathrm{grade}} &= mg\sin\theta
\end{align}

where
\[
\theta = \arctan(\mathrm{grade}).
\]

The total resistive force is
\[
F_{\mathrm{res}} = F_{\mathrm{aero}} + F_{\mathrm{roll}} + F_{\mathrm{grade}}.
\]

Using a spatial step \(ds\), the acceleration is computed from
\[
a = \frac{F_{\mathrm{drive}} - F_{\mathrm{res}}}{m},
\]
and the downstream speed is updated by the kinematic relation
\[
v_{i+1} = \sqrt{\max\left(0,\; v_i^2 + 2a\,ds\right)}.
\]

This is a simple but physically interpretable model. It captures the dominant terms that matter for a low-power competition vehicle and allows direct attribution of performance changes to specific physical causes.

\subsection{Powertrain Limits}

The model includes both a power limit and a force limit:
\[
F_{\mathrm{drive}}(v) = \min\left(F_{\max}, \frac{P_{\max}}{\max(v, v_\epsilon)}\right),
\]
where \(v_\epsilon\) is a small speed floor used to avoid singular behavior near zero speed. This form is important because low-speed launch behavior is often force-limited while moderate-speed cruising is power-limited.

\subsection{Cornering Constraint}

Track curvature \(\kappa(s)\) is converted to a local speed limit using a lateral acceleration cap:
\[
v_{\mathrm{corner}}(s) \le \sqrt{\frac{a_{y,\max}}{\kappa(s)}}.
\]

This constraint is central to race strategy. Even if a vehicle has enough power to accelerate, using that power before a tight corner can be wasteful if the controller must immediately force a coast or speed reduction before corner entry. The ability to account for curvature is therefore one of the main ways this framework improves over purely one-dimensional endurance models.

\subsection{Event Timing and Pace Requirement}

Let \(L_{\mathrm{event}}\) be the event distance and \(T_{\max}\) be the time limit. Then the average pace target is
\[
v_{\mathrm{avg,target}} = \frac{L_{\mathrm{event}}}{T_{\max}}.
\]

For the four-lap, \SI{35}{min} competition scenario used here,
\[
L_{\mathrm{event}} = 4 \times \SI{3850}{m} = \SI{15400}{m},
\qquad
T_{\max} = \SI{2100}{s},
\]
which yields
\[
v_{\mathrm{avg,target}} \approx \SI{7.33}{m/s} \approx \SI{26.4}{km/h}.
\]

In the present implementation, this target is not enforced through formal optimization; instead, it influences switching logic through a pace bias and burn/coast thresholds. This makes the controller lightweight and transparent, which is appropriate for early-stage strategy exploration and race-day interpretability.

\subsection{Energy Model}

Mechanical wheel power is
\[
P_{\mathrm{mech}} = F_{\mathrm{drive}}v.
\]

Fuel-cell electrical demand is approximated by
\[
P_{\mathrm{fc}} =
\begin{cases}
\dfrac{P_{\mathrm{mech}}}{\eta_{\mathrm{drive}}}, & P_{\mathrm{mech}} > 0 \\
0, & P_{\mathrm{mech}} = 0
\end{cases}
\]

and cumulative fuel-cell electrical energy is integrated over time:
\[
E = \int_0^T P_{\mathrm{fc}}(t)\,dt.
\]

Hydrogen mass is then estimated from fuel-cell efficiency \(\eta_{\mathrm{fc}}\) and hydrogen lower heating value \(LHV_{H_2}\):
\[
m_{H_2} = \frac{E}{\eta_{\mathrm{fc}}LHV_{H_2}}.
\]

This formulation is particularly useful for competition work because it bridges directly between simulation outputs and race metrics: elapsed time, electrical energy, and hydrogen mass.

It is also important to state what the model excludes. In the present vehicle architecture, the horn, onboard live telemetry system, and driver dashboard used for speed and strategy hints are powered by a separate accessories battery. Those loads are therefore not charged to the fuel-cell energy balance in this report.

\subsection{Rule-Based Burn-and-Coast Control}

The controller follows a rule-based logic:

\begin{itemize}[leftmargin=1.5em]
\item if the current speed is too low relative to the pace requirement, the controller burns;
\item if the current speed is sufficiently high, the controller coasts;
\item if lookahead detects an upcoming corner whose speed limit is below current speed plus a buffer, the controller transitions to coasting early to reduce entry losses.
\end{itemize}

This is not a globally optimal controller in the formal optimal-control sense, but it is well suited to the intended use case. It is transparent, tunable, and easy to connect to driver heuristics and engineering intuition.

\placeholderfigure
{Figure Placeholder 2: Speed, Corner Limit, and Burn/Coast Mode}
{2.5in}
{Recommended content: a distance-domain plot with actual speed, curvature-based corner speed limit, and shaded burn/coast regions.}
{Representative speed trace over distance. The ideal figure should show the actual vehicle speed, the local cornering speed limit, and highlighted regions where the controller applies propulsion or coasts. This plot is important because it makes the strategy interpretable: one can see whether energy is spent on useful acceleration or squandered before constrained segments.}

\section{Software Architecture}

\subsection{Track Processing Pipeline}

The first software layer converts raw measurement data into a reusable track model. The \texttt{export\_track\_from\_csv.py} module:

\begin{itemize}[leftmargin=1.5em]
\item detects latitude, longitude, and elevation columns,
\item filters invalid points,
\item converts geographic data to local distance increments,
\item resamples the path at fixed spatial resolution,
\item estimates grade from elevation change, and
\item estimates curvature from the change in heading angle.
\end{itemize}

This step is essential because raw GPS data is not immediately usable for strategy synthesis. Competition-relevant track modeling requires a smooth and spatially consistent representation of the path.

\subsection{Track Synchronization and Asset Generation}

The \texttt{sync\_tracks.py} utility automates the conversion of every CSV file in the track directory into a browser-ready JSON asset and generates a manifest used by the web dashboard. This is an important engineering decision: by standardizing the JSON schema, the same track data can be consumed by both the offline simulator and the browser interface.

\subsection{Simulation Engine}

The Python engine in \texttt{race\_strategy\_engine.py} performs the main offline analysis. It loads the selected track, tiles the track if necessary to reach the event distance, simulates burn-and-coast progression along the route, and exports both a full time series and a concise summary file. The simulation engine is the analytical core of the project.

\subsection{Interactive Dashboard}

The browser application performs the same conceptual tasks for interactive use:

\begin{itemize}[leftmargin=1.5em]
\item loads bundled tracks from the generated manifest,
\item accepts direct CSV upload for new tracks,
\item allows full-loop, start-point, and finish-point relocation by clicking on the track map,
\item reruns the simulation in the browser when parameters change, and
\item visualizes track geometry, speed, power, and cumulative energy.
\end{itemize}

The dashboard is valuable because it turns what would otherwise be a batch-only simulation into a practical engineering decision tool. Designers and strategy leads can vary mass, drag area, power limits, or controller thresholds and immediately observe the consequences.

\placeholderfigure
{Figure Placeholder 3: Dashboard Overview}
{2.6in}
{Recommended content: a screenshot of the web dashboard showing controls, KPIs, and the map. If possible, use a screenshot with a non-default start/finish selection to demonstrate track segmentation.}
{Interactive browser dashboard used for rapid scenario analysis. A complete figure should show the parameter controls, KPI cards, track map, speed gradient plots, and boundary-selection workflow. This figure demonstrates how the codebase supports engineering iteration without requiring direct edits to the simulation source.}

\section{Why This Matters for Shell Eco-marathon-Style Competition}

In competition, race strategy cannot be separated from design. A Shell Eco-marathon-style event rewards the vehicle that best balances two requirements: it must complete the route within the rules, and it must do so with minimal energy expenditure. This means the problem is inherently multidisciplinary.

The framework is useful to strategy creation in at least four ways:

\begin{enumerate}[leftmargin=1.5em]
\item \textbf{Track-specific planning.} The controller and plots identify where the vehicle should enter a burn phase, where it can safely coast, and where curvature or grade will dominate performance.
\item \textbf{Driver guidance.} Because the model is distance-based and visually tied to the track, it can be translated into driver cues such as accelerate out of this corner, coast through this sector, or preserve entry speed before a climb.
\item \textbf{Time-margin management.} Instead of treating energy minimization in isolation, the simulation continuously reveals whether the current setup is aggressive enough to maintain required average pace.
\item \textbf{Sensitivity analysis.} The same simulation environment can compare design variants before hardware is changed on the actual vehicle.
\end{enumerate}

The design implications are equally important. If the team tracks parameters such as \(m\), \(C_dA\), \(C_{rr}\), and cornering capability, the framework can show how each parameter changes event-level outcomes. This supports engineering prioritization:

\begin{itemize}[leftmargin=1.5em]
\item If aerodynamic drag dominates long straights, bodywork refinements may produce more value than a modest motor upgrade.
\item If rolling resistance dominates technical low-speed sections, tire and bearing choices may become more important than peak power.
\item If the accessories battery shows low state of charge during testing, the horn, live telemetry system, and driver dashboard should be treated as a separate electrical design problem rather than folded into the fuel-cell traction model.
\item If curvature repeatedly forces speed clipping, chassis balance and tire grip may unlock more competitive benefit than additional propulsion power.
\end{itemize}

This is why the project is more than a strategy calculator. It is a systems-engineering lens for evaluating where the next increment of effort should go.

\section{Example Results from the Current Project State}

Using the currently processed \texttt{sem\_2023\_us} track and the four-lap, \SI{35}{min} competition requirement, the example fuel-cell energy summary output is approximately:

\begin{itemize}[leftmargin=1.5em]
\item run time: \SI{1926.5}{s},
\item fuel-cell electrical energy: \SI{142.94}{Wh},
\item average speed: \SI{28.78}{km/h}, and
\item burn fraction: \(27.1\%\).
\end{itemize}

These values are not a final validation claim; rather, they provide a concrete example of the type of result the framework produces. The more important contribution is structural: the codebase now supports repeated regeneration of tracks from CSV, parameter sweeps, and map-aware boundary manipulation. The reported energy excludes accessories-battery loads by design.

\begin{table}[H]
\centering
\caption{Representative outputs produced by the current simulator configuration.}
\begin{tabular}{l c}
\toprule
\textbf{Metric} & \textbf{Value} \\
\midrule
Track name & \texttt{sem\_2023\_us} \\
Loop length & \SI{3850}{m} \\
Event distance (4 laps) & \SI{15400}{m} \\
Time limit & \SI{2100}{s} \\
Required average speed & \SI{26.40}{km/h} \\
Simulated run time & \SI{1926.52}{s} \\
Fuel-cell electrical energy & \SI{142.94}{Wh} \\
Simulated average speed & \SI{28.78}{km/h} \\
Burn fraction & 0.2710 \\
\bottomrule
\end{tabular}
\end{table}

\placeholderfigure
{Figure Placeholder 4: Energy and Power Demand}
{2.4in}
{Recommended content: paired plots of fuel-cell electrical power versus distance and cumulative traction energy versus distance, ideally for both a baseline and an improved design case.}
{Fuel-cell power demand and cumulative traction energy over the event distance. These plots are crucial for understanding whether improvements reduce only peak demand, only integrated energy, or both. They also help distinguish propulsion-related gains from reductions in aerodynamic or rolling losses, while excluding the separate accessories-battery loads from the fuel-cell accounting boundary.}

\section{Design Sensitivity and Parameter Tracking}

One of the strongest features of the project is that it can serve as a parameter-tracking framework. This is highly relevant to competition design reviews, where engineering teams must justify why a proposed modification is worth fabrication effort or test time.

The following parameter classes are especially important:

\subsection{Mass}

Mass influences acceleration, climbing losses, and rolling resistance. A lighter vehicle may exit corners more efficiently and reduce the energy penalty of repeated speed recovery.

\subsection{Aerodynamic Drag Area}

The term \(C_dA\) strongly affects power demand at speed. On circuits with longer straights or higher permissible cruise speed, drag reduction can have an outsized effect on hydrogen consumption.

\subsection{Rolling Resistance}

For low-speed endurance vehicles, rolling resistance remains important across nearly the entire lap. Tire pressure, compound, alignment, and bearing losses should therefore be viewed as race-strategy variables as much as mechanical details.

\subsection{Accessories Battery Boundary}

The accessories battery should still be tracked at the vehicle level, but it should not be conflated with traction energy. In the present architecture, the horn, onboard live telemetry system, and driver dashboard for speed and strategy hints are not powered by the fuel cell. That distinction matters because it prevents the traction model from overstating hydrogen consumption and keeps the fuel-cell analysis aligned with the actual electrical architecture.

\subsection{Cornering Capability}

Lateral acceleration capacity translates directly into allowable corner speed. Better cornering can reduce repeated deceleration and re-acceleration losses, which often matters more than marginal straight-line power gains on technical tracks.

\subsection{Propulsion Limits}

Motor force and fuel-cell power constraints shape where acceleration is possible and how quickly the vehicle can recover after constrained segments. Strategy and design are tightly coupled here: a vehicle with slightly better usable power may support a less aggressive pace threshold while still meeting timing requirements.

\section{Limitations and Recommended Future Work}

The present framework is intentionally lightweight and interpretable, but it includes several simplifying assumptions:

\begin{itemize}[leftmargin=1.5em]
\item constant air density and fixed drivetrain efficiency,
\item no wind model,
\item no thermal derating,
\item no battery or supercapacitor hybrid buffering,
\item no dynamic state-of-charge model for the separate accessories battery,
\item rule-based control instead of formal optimal control, and
\item limited validation against measured telemetry.
\end{itemize}

These limitations suggest a natural roadmap:

\begin{enumerate}[leftmargin=1.5em]
\item incorporate measured telemetry for calibration and residual analysis;
\item replace constant drive efficiency with a map-based efficiency model;
\item introduce stochastic disturbances such as wind and driver timing error;
\item add fuel-cell slew-rate and thermal constraints; and
\item benchmark the rule-based controller against dynamic programming or direct collocation.
\end{enumerate}

\section{Conclusion}

This project demonstrates that a technically rigorous race strategy tool can be built from a relatively compact but well-structured combination of physics, data processing, and software engineering. Its value for competition is twofold. First, it improves race strategy by locating where energy should and should not be spent, with explicit awareness of timing requirements and track geometry. Second, it helps the team understand how design decisions affect competition performance by connecting measurable parameters such as mass, drag area, rolling resistance, and cornering capability to event-level outcomes.

For a competition like Shell Eco-marathon, that connection is essential. Vehicles do not win solely by being low-drag or light in the abstract; they win by converting those properties into better event performance on a real track under real constraints. A framework that unifies strategy and design analysis is therefore not just useful, but strategically important.

\appendix

\section{Representative Implementation Snippets}

This appendix includes short excerpts showing where the main capabilities are implemented in the codebase. These snippets are not intended as complete source listings; instead, they highlight the project components most relevant to the modeling and workflow described in the main text.

\subsection{Track CSV Processing and Resampling}

\begin{lstlisting}[style=projectcode, language=Python, caption={Representative logic for converting raw CSV track data into a simulation-ready JSON representation.}]
def build_track_from_dataframe(df, track_name, ds_grid=2.0):
    lat_col = find_column(df, ['Latitude', 'Lat'])
    lon_col = find_column(df, ['Longitude', 'Lon', 'Lng', 'Long'])
    alt_col = find_column(
        df,
        ['Metres above sea level', 'Meters above sea level',
         'Elevation', 'Altitude', 'Alt'],
        required=False,
    )

    lat = pd.to_numeric(df[lat_col], errors='coerce').to_numpy(float)
    lon = pd.to_numeric(df[lon_col], errors='coerce').to_numpy(float)
    alt = np.zeros_like(lat) if alt_col is None else \
        pd.to_numeric(df[alt_col], errors='coerce').ffill().bfill().to_numpy(float)

    R = 6371000.0
    lat_r = np.deg2rad(lat)
    lon_r = np.deg2rad(lon)
    dlat = np.diff(lat_r, prepend=lat_r[0])
    dlon = np.diff(lon_r, prepend=lon_r[0])
    dy = R * dlat
    dx = R * np.cos(lat_r) * dlon
    ds = np.sqrt(dx**2 + dy**2)
    s_raw = np.cumsum(ds)
    s_raw[0] = 0.0

    s = np.arange(0.0, s_raw[-1] + ds_grid, ds_grid)
    lat_g = np.interp(s, s_raw, lat)
    lon_g = np.interp(s, s_raw, lon)
    alt_g = np.interp(s, s_raw, alt)
\end{lstlisting}

\subsection{Longitudinal Simulation and Burn-and-Coast Logic}

\begin{lstlisting}[style=projectcode, language=Python, caption={Core Python simulation loop used for the offline race strategy engine.}]
for i in range(n - 1):
    dsi = max(ds[i + 1], 1e-3)
    look_idx = min(n - 1, i + max(1, int(sp.corner_lookahead_m / dsi)))
    future_corner_limit = np.min(v_corner[i:look_idx + 1])

    if v[i] > future_corner_limit + sp.corner_buffer_mps:
        burning = False
    elif v[i] < min(sp.burn_on_threshold_mps, target_avg + 0.2):
        burning = True
    elif v[i] > max(sp.burn_off_threshold_mps, target_avg + 0.8):
        burning = False

    f_res = resistive_force(v[i], track['grade'][i], p)
    if burning:
        drive_power = min(p.motor_power_max_w, p.fuel_cell_peak_usable_power_w)
        f_drv = min(p.motor_force_max_n, drive_power / max(v[i], 0.5))
    else:
        f_drv = 0.0

    a = np.clip((f_drv - f_res) / p.mass_kg,
                -p.brake_max_mps2, p.accel_max_mps2)
    v_next = np.sqrt(max(0.0, v[i] ** 2 + 2 * a * dsi))
    v_next = min(v_next, v_corner[i + 1], p.abs_speed_cap_mps)
\end{lstlisting}

\subsection{Browser-Based Start and Finish Relocation}

\begin{lstlisting}[style=projectcode, language=JavaScript, caption={Interactive start/finish relocation logic in the web dashboard.}]
mapDiv.on('plotly_click', (event) => {
  if (!boundarySelection.pendingSelection) {
    return;
  }

  const clickedPoint = event.points?.find(
    point => point?.curveNumber === 2 || point?.pointNumber != null
  );
  const clickedIndex = clickedPoint?.pointNumber;
  if (!Number.isInteger(clickedIndex)) {
    return;
  }

  const wasFullLoop = boundarySelection.startIndex === boundarySelection.finishIndex;
  if (boundarySelection.pendingSelection === 'start') {
    boundarySelection.startIndex = clickedIndex;
    if (wasFullLoop) {
      boundarySelection.finishIndex = clickedIndex;
    }
  } else if (boundarySelection.pendingSelection === 'finish') {
    boundarySelection.finishIndex = clickedIndex;
  }

  boundarySelection.pendingSelection = null;
  updateBoundaryControls(baseTrack);
  runScenario();
});
\end{lstlisting}

\subsection{Track Synchronization for Browser Assets}

\begin{lstlisting}[style=projectcode, language=Python, caption={Utility that exports all CSV tracks and builds the browser manifest.}]
for csv_path in sorted(CSV_TRACK_DIR.glob('*.csv')):
    out_path = DOCS_TRACK_DIR / f'{csv_path.stem}.json'
    track = export_track(str(csv_path), str(out_path), ds_grid=ds_grid)
    manifest.append({
        'label': prettify_name(csv_path.stem),
        'path': f'tracks/{out_path.name}',
        'kind': 'generated',
        'source_csv': csv_path.name,
        'length_m': round(float(track['length_m']), 3),
        'point_count': len(track['distance_m']),
    })
\end{lstlisting}

\end{document}
